% Generated draft: requires mathematical and source review.
\documentclass[11pt]{article}
\usepackage[margin=25mm]{geometry}
\usepackage{fontspec}
\setmainfont{Latin Modern Roman}
\usepackage{amsmath,amssymb,mathtools,unicode-math}
\setmathfont{Latin Modern Math}
\usepackage{enumitem,needspace,xurl,hyperref}
\setlength{\parindent}{0pt}
\setlength{\parskip}{5pt}
\setlength{\emergencystretch}{3em}
\newcommand{\N}{\mathbb{N}}
\newcommand{\Z}{\mathbb{Z}}
\newcommand{\Q}{\mathbb{Q}}
\newcommand{\R}{\mathbb{R}}
\newcommand{\C}{\mathbb{C}}
\newcommand{\F}{\mathbb{F}}
\newcommand{\A}{\mathbb{A}}
\newcommand{\PP}{\mathbb{P}}
\newcommand{\E}{\mathbb{E}}
\newcommand{\eps}{\varepsilon}
\newcommand{\dd}{\,\mathrm{d}}
\begin{document}
\section*{UnsolvedMath ID: 30000959}
\textbf{Source problem code:} catalogue entry
\subsection{Definitions and mathematical statement}
Let \(G\) be a complex reductive algebraic group. For a finitely generated group \(\Gamma\), use the customary affine character variety
\[
X_G(\Gamma)=\operatorname{Hom}(\Gamma,G)\mathbin{/\!/}G,
\]
where the quotient is by conjugation. For a knot \(K\subset S^3\), let \(M_K=S^3\setminus\operatorname{int}\nu K\) be its exterior, with boundary torus \(T_K\). Its fundamental group identifies with that of \(S^3\setminus K\). Inclusion of the boundary induces the restriction morphism
\[
r_K:X_G(\pi_1M_K)\longrightarrow X_G(\pi_1T_K).
\]
After a peripheral identification \(\pi_1T_K\cong\Z^2\), write \(\mathcal I_G(K)=\operatorname{im}(r_K)\subseteq X_G(\Z^2)\).

Conjecture 7 in the source asks whether the quantum \(G\)-invariants determine this boundary image. If \(Q_G(K)\) denotes the full collection of quantum invariants intended there, the stated equivalent implication is
\[
\forall G\ \forall K_1,K_2\subset S^3,\qquad
Q_G(K_1)=Q_G(K_2)\quad\Longrightarrow\quad
\mathcal I_G(K_1)=\mathcal I_G(K_2).
\]
Comparison of the images requires compatible peripheral identifications, customarily using meridian and longitude. The source does not fix their orientation or framing conventions. It also does not specify the family of quantum invariants, allowed colors or representations, quantum parameters, or normalizations. These omissions prevent interpreting equality of \(Q_G\) as one particular known knot invariant.

The displayed image is the literal image requested in the extraction. Replacing it by its Zariski closure or a scheme-theoretic image is a different formulation and is not silently done. Likewise, the use of the customary affine character variety above does not resolve any source-specific reduced or scheme-level convention.
\subsection{Short English statement}
For each complex reductive group, must knots with the same intended quantum invariants have the same image of their character varieties in the boundary-torus character variety? The quantum family, peripheral identifications, and meaning of image need source review.
\subsection{Sources}
\begin{itemize}
\item Local UnsolvedMath registry, permanent ID \(30000959\), source code catalogue entry, ``Quantum Knot Invariants and Boundary Character Varieties.'' The original labels the question Conjecture 7; no individual proposer is supplied.
\item \emph{Invariants in Low-Dimensional Topology} (2009), Oberwolfach Reports, DOI 10.4171/owr/2008/22, pp. 1157--1232: \url{https://doi.org/10.4171/owr/2008/22}.
\item Supplied reference, ``Sikora 2008'': \url{https://arxiv.org/abs/0807.0943}.
\item Supplied literature triage, checked 2026-08-21: status ``open.'' It calls this an instance of the quantum-character-variety determination problem and reports no found general resolution. This dated assessment derives from supplied OWR curation; time-sensitive status and cited claims require independent verification.
\end{itemize}
\end{document}
